% Part: sets-functions-relations
% Chapter: functions
% Section: kinds

\documentclass[../../../include/open-logic-section]{subfiles}

\begin{document}

\olfileid{sfr}{fun}{kin}
\olsection{Soorten functies}

\begin{explain}
Het is nuttig een indeling te geven van enkele soorten functies die we
het vaakst tegenkomen.

Om te beginnen bekijken we functies met de eigenschap dat ieder
element van het codomein een waarde van de functie is. Zulke functies
heten !!{surjective} en kunnen worden weergegeven zoals in
\olref{fig:surjective}.

\begin{figure}
  \olasset{assets/diagrams/surjective.tikz}
  \caption{!!^a{surjective} functie heeft ieder !!{element} van het
    codomein als waarde.}
  \ollabel{fig:surjective}
\end{figure}
\end{explain}

\begin{defn}[!!^{surjective} functie]
Een functie $f \colon A \rightarrow B$ heet \emph{!!{surjective}} dan
en slechts dan als $B$ ook het bereik van~$f$ is, dat wil zeggen: voor
iedere $y \in B$ bestaat er ten minste één $x \in A$ waarvoor~$f(x) =
y$, of in symbolen:
\[
  (\forall y \in B)(\exists x \in A)f(x) = y.
\]
Zo'n functie noemen we !!a{surjection} van $A$ naar $B$.
\end{defn}

\begin{explain}
Om aan te tonen dat $f$ !!a{surjection} is, moet men laten zien dat
ieder object in het codomein van $f$ de waarde $f(x)$ is voor een
zekere invoer $x$.

Merk op dat iedere functie !!a{surjection} \emph{induceert}. Gegeven
een functie $f \colon A \to B$ definiëren we immers $f' \colon A \to
\ran{f}$ door $f'(x) = f(x)$. Omdat $\ran{f}$ \emph{per definitie}
gelijk is aan $\Setabs{f(x) \in B}{x \in A}$, is deze functie $f'$
gegarandeerd !!a{surjection}.
\end{explain}

\begin{explain}
Iedere functie beeldt elke mogelijke invoer af op precies één uitvoer.
Er zijn echter ook functies die nooit verschillende invoeren op
dezelfde uitvoer afbeelden. Zulke functies heten !!{injective} en
kunnen worden weergegeven zoals in \olref{fig:injective}.
\begin{figure}
  \olasset{assets/diagrams/injective.tikz}
  \caption{!!^a{injective} functie beeldt nooit twee verschillende
    argumenten af op dezelfde waarde.}
  \ollabel{fig:injective}
\end{figure}
\end{explain}

\begin{defn}[!!^{injective} functie]
Een functie $f \colon A \rightarrow B$ heet \emph{!!{injective}} dan
en slechts dan als er voor iedere $y \in B$ ten hoogste één $x \in A$
bestaat waarvoor~$f(x) = y$. Zo'n functie noemen we !!a{injection} van
$A$ naar~$B$.
\end{defn}

\begin{explain}
Om aan te tonen dat $f$ !!a{injection} is, moet men voor willekeurige
!!{element}s $x$ en $y$ van het domein van $f$ bewijzen dat uit
$f(x)=f(y)$ volgt dat $x=y$.
\end{explain}

\begin{ex}
De constante functie $f\colon \Nat \to \Nat$ gegeven door $f(x) = 1$
is noch !!{injective}, noch !!{surjective}.

De identiteitsfunctie $f\colon \Nat \to \Nat$ gegeven door $f(x) = x$
is zowel !!{injective} als !!{surjective}.

De opvolgerfunctie $f \colon \Nat \to \Nat$ gegeven door $f(x) = x+1$
is !!{injective}, maar niet !!{surjective}.
  
De functie $f \colon \Nat \to \Nat$ gedefinieerd door
\[
  f(x) =
  \begin{cases}
    \frac{x}{2} & \text{als $x$ even is} \\
    \frac{x+1}{2} & \text{als $x$ oneven is.}
  \end{cases}
\]
is !!{surjective}, maar niet !!{injective}.
\end{ex}

\begin{explain}
Vaak willen we functies beschouwen die zowel !!{injective} als
!!{surjective} zijn. Zulke functies noemen we !!{bijective}. Ze zien
eruit als de functie in \olref{fig:bijective}. !!^{bijection}s worden
ook wel \emph{een-eenduidige correspondenties} genoemd, omdat ze de
elementen van het codomein eenduidig koppelen aan de elementen van het
domein.
\begin{figure}
  \olasset{assets/diagrams/bijective.tikz}
  \caption{!!^a{bijective} functie koppelt de elementen van het
    codomein eenduidig aan die van het domein.}
  \ollabel{fig:bijective}
\end{figure}
\end{explain}

\begin{defn}[!!^{bijection}]
Een functie $f \colon A \to B$ heet \emph{!!{bijective}} dan en
slechts dan als zij zowel !!{surjective} als !!{injective} is. Zo'n
functie noemen we !!a{bijection} van $A$ naar~$B$ (of tussen $A$
en~$B$).
\end{defn}
\end{document}
